\section{Methodology}
\label{sec:methodology}

Describe \emph{how} you carried out the work. This section should be detailed
enough for another researcher to replicate the study.

\subsection{Data / Materials}

Describe the data sources, datasets, or materials used.
\Cref{tab:dataset} summarizes the dataset characteristics.

\begin{table}[h]
  \centering
  \caption{Dataset overview.}
  \label{tab:dataset}
  \begin{tabular}{lrr}
    \toprule
    Split      & Samples & Features \\
    \midrule
    Training   & 8\,000  & 128      \\
    Validation & 1\,000  & 128      \\
    Test       & 1\,000  & 128      \\
    \midrule
    \textbf{Total} & \textbf{10\,000} & 128 \\
    \bottomrule
  \end{tabular}
\end{table}

\subsection{Approach}

Describe the method, algorithm, or process step by step.
\Cref{fig:pipeline} illustrates the processing pipeline.

\begin{figure}[h]
  \centering
  \begin{tikzpicture}[
    node distance=2.2cm,
    stage/.style={rectangle, draw, rounded corners=4pt,
                  minimum width=2.4cm, minimum height=0.9cm,
                  fill=codeblue!12, font=\small},
  ]
    \node[stage] (a) {Data};
    \node[stage] (b) [right of=a] {Process};
    \node[stage] (c) [right of=b] {Evaluate};
    \node[stage] (d) [right of=c] {Output};

    \draw[->, thick] (a) -- (b);
    \draw[->, thick] (b) -- (c);
    \draw[->, thick] (c) -- (d);
  \end{tikzpicture}
  \caption{Overview of the processing pipeline.}
  \label{fig:pipeline}
\end{figure}

\subsection{Evaluation Metrics}

Define the metrics used to measure success:

\begin{description}
  \item[Metric A] Formal definition, e.g.\ $F_1 = 2\,\frac{P \cdot R}{P + R}$
  \item[Metric B] Another metric with its formula or description
\end{description}
